% librerie matematica
\usepackage{siunitx} % fornisce dei typeset per i numeri e le unita di misura
\usepackage{amssymb}
\usepackage{amsthm}

\usepackage{amsfonts}

% emoji
%\usepackage{fourier}
% \usepackage{emoji}  % richiede LuaLaTeX + font emoji a colori

% creazione dei Toast message
\usepackage[many]{tcolorbox}

% definizione di nuovi colori
\usepackage{xcolor}

% SIMBOLO MATEMATICO ESISTE ED E' UNICO
\newcommand{\existsone}{\exists !}

% INSIEMI MATEMATICI
\newcommand{\R}{\mathbb{R}}
\newcommand{\N}{\mathbb{N}}
\newcommand{\Z}{\mathbb{Z}}
\newcommand{\C}{\mathbb{C}}

% DEFINIZIONE
\definecolor{myred}{RGB}{191,101,127}
\definecolor{mygreen}{RGB}{0,168,107}

\definecolor{myblue}{RGB}{0,127,255}
\definecolor{myblue2}{RGB}{36,84,255}

\definecolor{myyellow}{RGB}{255,191,0}
\definecolor{mygray}{RGB}{83,104,120}

% BOX VUOTI
\newtcolorbox{empty_green}{enhanced, toprule=0pt, bottomrule=0pt, rightrule=0pt, leftrule=3pt, arc=0mm, skin=enhancedlast jigsaw, sharp corners,
colframe=mygreen, colbacktitle=mygreen}

\newtcolorbox{empty_red}{enhanced, toprule=0pt, bottomrule=0pt, rightrule=0pt, leftrule=3pt, arc=0mm, skin=enhancedlast jigsaw, sharp corners,
colframe=myred, colbacktitle=myred}

\newtcolorbox{empty_blue}{enhanced, toprule=0pt, bottomrule=0pt, rightrule=0pt, leftrule=3pt, arc=0mm, skin=enhancedlast jigsaw, sharp corners,
colframe=myblue, colbacktitle=myblue}

\newtcolorbox{empty_yellow}{enhanced, toprule=0pt, bottomrule=0pt, rightrule=0pt, leftrule=3pt, arc=0mm, skin=enhancedlast jigsaw, sharp corners,
colframe=myyellow, colbacktitle=myyellow}


\newcommand{\BoxRed}[3]{
    \begin{empty_red}
        \textcolor{myred}{\textbf{#1} \textbf{#2}}\\
        \noindent #3
    \end{empty_red}
}

\newcommand{\BoxGreen}[3]{
    \begin{empty_green}
        \textcolor{mygreen}{\textbf{#1} \textbf{#2}}\\
        \noindent #3
    \end{empty_green}
}

\newcommand{\BoxBlue}[3]{
    \begin{empty_blue}
        \textcolor{myblue}{\textbf{#1} \textbf{#2}}\\
        \noindent #3
    \end{empty_blue}
}

\newcommand{\Attenzione}[1]{
    \begin{empty_yellow}
        \textbf{\Large !}\enspace #1
    \end{empty_yellow}
}

\newcommand{\Definizione}[2]{
    \begin{empty_green}
         \textcolor{mygreen}{\textbf{Definizione (#1)}}\\
        \noindent #2
    \end{empty_green}
}

\newcommand{\Teorema}[2]{
    \begin{empty_red}
        \textcolor{myred}{\textbf{Teorema (#1)}}\\
        \noindent #2
    \end{empty_red}
}

\newcommand{\Dimostrazione}[1]{
    \begin{empty_blue}
        \textcolor{myblue}{\textbf{Dimostrazione}}\\
        \noindent \textcolor{myblue}{#1}
    \end{empty_blue}
}